\documentclass[openany]{book}

\usepackage[margin=1in]{geometry}
\usepackage{amsmath,amsfonts,amsthm, amssymb}
\usepackage{yhmath}
\usepackage{mathrsfs}
\usepackage{mathtools}
\usepackage{xcolor}
\usepackage{graphicx}
\usepackage{comment}
\usepackage{tikz-cd}
\usepackage{quiver}
\usepackage{hyperref,cleveref}
\renewcommand{\familydefault}{ppl}
\newcommand{\tr}{\text{tr}}
\newcommand{\R}{\mathbb{R}}
\newcommand{\E}{\mathbb{E}}
\newcommand{\Z}{\mathbb{Z}}
\newcommand{\C}{\mathbb{C}}
\newcommand{\F}{\mathbb{F}}
\newcommand{\la}{\langle}
\newcommand{\ra}{\rangle}
\newcommand{\colim}{\text{colim}}
\DeclareMathOperator{\im}{im}
\DeclareMathOperator{\disc}{disc}
\let\oldemptyset\emptyset
\let\emptyset\varnothing
\newcommand{\tor}{\text{Tor}}
\newcommand{\id}{\text{id}}
\newcommand{\ext}{\text{Ext}}
\newcommand{\ptop}{\text{PTop}}
\newcommand{\pt}{\text{pt}}
\newcommand{\ach}{\text{Ach}}
\newcommand{\Q}{\mathbb{Q}}
\newcommand{\gal}{\text{Gal}}
\newcommand{\fraccomma}{\genfrac{}{}{0pt}{}{}{,}}
\newcommand{\diverg}{\operatorname{div}}
\newcommand{\curl}{\operatorname{curl}}
\definecolor{wikipediadarkblue}{rgb}{0.023, 0.270, 0.676}
\hypersetup{
    colorlinks,
    citecolor=black,
    filecolor=black,
    linkcolor=wikipediadarkblue,
    urlcolor=red
}
\usepackage{pgfplots}
\pgfplotsset{compat=newest}


\input{hui_r.tex}

\title{Calc II Section Notes
\\ 
\vspace{0.4cm}
\large Fall 2026}




\date{\today}
\author{Hui Sun\\
(hsun95@jh.edu)}


\begin{document}

\maketitle
\tableofcontents

% \tableofcontents
\newpage

\chapter{Applications of Integrals}


\section*{\centering Week 1 (8/31-9/3)}



\begin{itemize}
    \item TA: Hui.
    \item Email: hsun95@jh.edu.
    \item Office Hours: Wednesday 1-2 PM, Krieger 211.
\end{itemize}




\begin{prob}[discuss with your teammates]
    You are an open water diver, which of the following sharks would you be excited to see? Rank from 1 to 5 most excited to least excited. 
\end{prob}

% \begin{center}
%     \begin{tabular}{|l|p{7cm}|c|}
%     \hline
%     \textbf{Shark Type} & \textbf{Diet} &
%     \textbf{Confirmed human attacks} \\ \hline
    
%     Tiger shark &
%     Sea turtles, rays, other sharks, fish, seabirds, dolphins,
%     squid, crustaceans &
%     142 (39 fatal)\\ \hline
    
%     Hammerhead shark &
%     Fish, squid, crustaceans, rays, and smaller sharks &
%     18 (0 fatal)\\ \hline
    
%     Nurse shark &
%     Fish, stingrays, octopus, squid, clams, and crustaceans &
%     9 (0 fatal)\\ \hline
    
%     Thresher shark &
%     Mostly schooling fish such as herring and mackerel; also squid &
%     0 listed \\ \hline
    
%     Great white shark &
%     Fish, rays, other sharks, squid, and marine mammals &
%     351 (59 fatal)\\ \hline
    
%     \end{tabular}
% \end{center}


\begin{center}
    \begin{tabular}{|l|p{5cm}|c|p{6cm}|}
    \hline
    \textbf{Shark Type} &
    \textbf{Diet} &
    \textbf{Humans} &
    \textbf{Appearance} \\ \hline
    
    Tiger shark &
    Sea turtles, rays, other sharks, fish, seabirds, dolphins,
    squid, crustaceans, and carrion &
    142 (39 fatal)& Long body with a broad head and dark vertical bands
    along the sides.
     \\ \hline
    
    Hammerhead shark &
    Fish, squid, crustaceans, rays, and smaller sharks &
    18 (0 fatal) &
    Distinctive hammer-shaped head with eyes positioned far apart.
    \\ \hline
    
    Nurse shark &
    Fish, stingrays, octopus, squid, clams, and crustaceans &
    9 (0 fatal)&
    Rounded head with small eyes and two barbels near the mouth.
    \\ \hline
    
    Thresher shark &
    Mostly schooling fish such as herring and mackerel; also squid &
    0 listed &
    Extremely long upper tail fin,
    nearly as long as the rest of the body. \\ \hline
    
    Great white shark &
    Fish, rays, other sharks, squid, and marine mammals &
    351 (59 fatal)&
    Muscular body with prominent
    triangular teeth. Powerful and heavy-looking.\\ \hline
    
    \end{tabular}
    \end{center}


\begin{prob}
    Approximate the area under $f(x)=x^2$
    on $[0,1]$ using Riemann sums with $N=2,3$
    and some choice of sample points. Then think about the limit as $N\to\infty$. Hint: 
    \begin{equation*}
        \sum_{i=1}^N i=1+2+\dots+N=\frac{N(N+1)}{2}, \quad \sum_{i=1}^N i^2=1+4+9+\dots+N^2=\frac{N(N+1)(2N+1)}{6}
    \end{equation*}
\end{prob}
\begin{proof}
    We will use left endpoints. When $N=2$, $x_1=0, x_2=1/2$. Thus 
    \begin{equation*}
        R_2=0+\frac{1}{4}\cdot\frac{1}{2}=\frac{1}{8}
    \end{equation*}
    When $N=3$, $x_1=0, x_2=1/3, x_3=2/3$. Then 
    \begin{equation*}
        R_3=0+\frac{1}{9}\cdot\frac{1}{3}+\frac{4}{9}\cdot\frac{1}{3}=\frac{5}{27}
    \end{equation*}
    Now we compute when $N\to\infty$ 
    \begin{equation*}
        R=\lim_{N\to\infty}\sum_{i=0}^{N-1}\left(\frac{i}{N}\right)^2\frac{1}{N}=\frac{1}{6}
    \end{equation*}
\end{proof}



\begin{prob}
    Compute and interpret the limit
    \[
    \lim_{N\to\infty}
    \sum_{i=1}^{N}
    \left(
        \frac{i}{N^2}
        -
        \frac{i^2}{N^3}
    \right)
    \]
\end{prob}
\begin{proof}
     We first compute it:
    \begin{align}
        \sum_{i=1}^{N}
    \left(
        \frac{i}{N^2}
        -
        \frac{i^2}{N^3}
    \right)&=\frac{N(N+1)}{2N^2}-\frac{N(N+1)(2N+1)}{6N^3}
    \end{align}
    Letting $N\to\infty$, 
    \begin{equation*}
        \lim_{N\to\infty}
    \sum_{i=1}^{N}
    \left(
        \frac{i}{N^2}
        -
        \frac{i^2}{N^3}
    \right)=\frac{1}{2}-\frac{1}{3}=\frac{1}{6}
    \end{equation*}
    Now to interpret it, we can write it as a Riemann sum:
    \begin{equation*}
        \sum_{i=1}^{N}
    \left(
        \frac{i}{N^2}
        -
        \frac{i^2}{N^3}
    \right)=\sum_{i=1}^N\frac{1}{N}\left(\frac{i}{N}-\frac{i^2}{N^2}\right)
    \end{equation*}
    Then $\Delta x=\frac{1}{N}$, and $x_i=\frac{i}{N}$. Letting $N\to\infty$, this is the Riemann sum of the following
    \begin{equation*}
        \int_0^1 (x-x^2)\,dx
    \end{equation*}
\end{proof}


\begin{prob}
    Compute
    \[
    \int_{-1}^{5} (1+3x)\,dx
    \]
    using Riemann sums.
\end{prob}
\begin{proof}
    We divide the interval $[-1, 5]$ into $N$ intervals, then $\Delta x=6/N$, and if we use right endpoints, the $i$th sample point is $x_i=-1+6i/N$. Then 
    \begin{align*}
        \sum_{i=1}^Nf(x_i)\Delta x&=\sum_{i=1}^N(-2+18i/N)\frac{6}{N}
    \end{align*}
    letting $N\to\infty$, the answer is $42$.
\end{proof}




% \begin{prob}
%     Let $f$ be an increasing function on $[0,1]$. Is it true that the Riemann sum using left endpoints is less than $\int_0^1 f(x)dx$?
% \end{prob}
% \begin{proof}
%     It is true.
% \end{proof}



\newpage
\section*{\centering Week 2 (9/7-9/11)}
\begin{prob}
    Compute the following integral:
    \begin{equation*}
        \int_1^3\frac{\ln x}{x(1+\ln^2x)}dx
    \end{equation*}
\end{prob}
\begin{proof}
    Let $u=1+\ln^2 x$, then $du=\frac{2\ln x}{x}dx$, and 
    \begin{align*}
        \int_1^3\frac{\ln x}{x(1+\ln^2x)}dx&=\int_{1}^{1+\ln^23}\frac{1}{2u}du=\frac{1}{2}\ln(1+\ln^23)
    \end{align*}
\end{proof}


\begin{prob}
    Compute the following integral:
    \begin{equation*}
        \int_0^{1/2}\frac{x}{\sqrt{1-x^4}}dx
    \end{equation*}
\end{prob}
\begin{proof}
    Let $u=x^2, du=2xdx$, then 
    \begin{equation*}
        \int_0^{1/2}\frac{x}{\sqrt{1-x^4}}dx=\int_0^{1/4}\frac{1}{2\sqrt{1-u^2}}du=\frac{1}{2}\left(\arcsin 1/4-\arcsin 0\right)
    \end{equation*}
\end{proof}

\begin{prob}
    Compute the following integral:
    \begin{equation*}
        \int_1^2\frac{e^x}{(1+e^x)\ln(1+e^x)}dx
    \end{equation*}
\end{prob}
\begin{proof}
    Let $u=\ln(1+e^x)$, then $du=\frac{e^x}{1+e^x}dx$, thus 
    \begin{equation*}
        \int_1^2\frac{e^x}{(1+e^x)\ln(1+e^x)}dx=\int_{\ln(1+e)}^{\ln(1+e^2)}\frac{1}{u}du=\ln\left(\frac{\ln(1+e^2)}{\ln(1+e)}\right)
    \end{equation*}
\end{proof}


\newpage
\section*{\centering Week 3 (9/14-9/18)}
% Ball example.

% For the washer method, if rotated about the $x$-axis, integrate wrt $x$; if rotated about the $y$-axis, integrate wrt $y$.
\begin{prob}
    Compute the area enclosed by the following curves:
    \begin{enumerate}
        \item \begin{equation*}
            y=\frac{1}{x}, y=\frac{1}{x^2}, x=2
        \end{equation*}
        \item \begin{equation*}
            x=1-y^2, x=y^2-1
        \end{equation*}
        \item \begin{equation*}
            x=\frac{\pi}{2}, x=\pi, y=\sin x, y=x
        \end{equation*}
        \item \begin{equation*}
            y=\sqrt[3]{2x}, y=\frac{1}{8}x^2, x=0, x=6
        \end{equation*}
    \end{enumerate}
\end{prob}
\begin{proof}
    \begin{enumerate}
        \item \begin{equation*}
            \int_1^2\left(\frac{1}{x}-\frac{1}{x^2}\right)dx
        \end{equation*}
        \item \begin{equation*}
            \int_{-1}^1(1-y^2)-(y^2-1)dy
        \end{equation*}
        \item \begin{equation*}
            \int_{\pi/2}^\pi\left(x-\sin x\right)dx
        \end{equation*}
        \item \begin{equation*}
            \int_0^4\left(\sqrt[3]{2x}-x^2/8\right)dx+\int_4^6\left(x^2/8-\sqrt[3]{2x}\right)dx
        \end{equation*}
    \end{enumerate}
\end{proof}


\begin{prob}
    Write down the integrals for the following solids of revolution using the washer method. DO NOT SOLVE.
    \begin{enumerate}
        \item $y=x, y=x^3, x\geq 0$, rotated about the $x$-axis.
        \item $y^2=x, x^2=y$, rotated about $y=1$.
        \item $y^2=x, x^2=y$, rotated about the $y$-axis.
        \item $x=0, y=x^2, y=-2x+3$ around the $y$-axis.
    \end{enumerate}
\end{prob}
\begin{proof}
    \begin{enumerate}
        \item The cross-sectional area in the $x$-direction is $\pi(x^2-(x^3)^2)$, thus 
        \begin{equation*}
            S=\int_0^{1}\pi(x^2-x^6)dx
        \end{equation*}
        \item The cross-sectional area along the $x$-axis is 
        \begin{equation*}
            A=\pi\left[(1-x^2)^2-(1-\sqrt{x})^2\right]
        \end{equation*}
        Thus the solid volume is 
        \begin{equation*}
            S=\int_0^1\pi\left[(1-x^2)^2-(1-\sqrt{x})^2\right]dx
        \end{equation*}
        \item The cross-sectional area along the $y$-axis is 
        \begin{equation*}
            \pi\left[(\sqrt{y})^2-(y^2)^2\right]dy
        \end{equation*}
        Thus the solid volume is 
        \begin{equation*}
            S=\int_0^1\pi\left[(\sqrt{y})^2-(y^2)^2\right]dy
        \end{equation*}
        \item This needs to be separated into two terms:
        \begin{equation*}
            S=\int_0^1\pi(\sqrt{y})^2dy+\int_1^3\pi\left(\frac{(3-y)}{2}\right)^2dy
        \end{equation*}
    \end{enumerate}
\end{proof}

% The solid obtained by rotating the disk of radius 1 centered at $(2,0)$ in the $xy$-plane around the $y$-axis.


% \textbf{Solutions:} We can use the cylindrical shell method for both.
% \begin{enumerate}
%     \item Setting $x^2=-2x+3$, we see these two curves intersect at $x=1$ in the first quadrant. One can draw a picture to see $-2x+3\geq x^2$ on $[0,1]$. Thus, 
%     \begin{equation}
%         \text{Vol}=\int_0^1 (-2x+3-x^2)\cdot (2\pi x)dx
%     \end{equation}
%     \item For a fixed $x\in [1,3]$, we see the ``height'' is given by $2\sqrt{1-(x-2)^2}$, thus 
%     \begin{equation*}
%         \text{Vol}=\int_1^3 2\sqrt{1-(x-2)^2}\cdot (2\pi x)dx
%     \end{equation*}
% \end{enumerate}




\end{document}